% ============================================================
%  frontmatter/abstract.tex — Dissertation Abstract
% ============================================================

\chapter*{Abstract}
\addcontentsline{toc}{chapter}{Abstract}

\textit{150-250 words summary}

The rapid deployment of Large Language Model (LLM) powered infrastructure in higher education is driven primarily by efficiency gains. However, this unchecked integration presents a significant challenge for institutions dedicated to nurturing critical thinking, diverse perspectives, and intellectual empowerment. By inherently privileging statistically frequent patterns, LLMs risk functioning as ``standardizing engines'' that compress diverse explanatory styles and localized knowledge systems into a narrow, homogenized norm—a phenomenon conceptualized in educational theory as Epistemic Enclosure. 

This study audits the extent to which LLM-generated mentoring homogenizes pedagogical explanatory styles and intellectual perspectives in Graduate Teaching Assistant (GTA) workflows. By decoupling upstream machine-generated artifacts (Dimension A) from downstream user cognitive agency (Dimension B), this research implements a task-anchored computational audit comparing routine concept explanations (Task 1: objective) with open-ended project scoping (Task 2: subjective). Utilizing zero-shot API execution across commercial foundation models (GPT-4o, Claude 3.5 Sonnet, and Gemini 1.5 Pro) combined with high-dimensional vector metrics (Sentence Embedding Variance and BERTopic clustering), this study establishes an objective, scalable system benchmark bridging data science with critical educational theory.

\newpage